\section{Seed Corpus and Dictionary}
We have chosen to use the seed set 'pngSuite' from Willem van Schaik \cite{vanschaikPngSuiteOfficialSet} in order to have a wide variety of PNGs (with different color-type, bit-depth and interlacing).
The seeds are preferably small images to minimize execution time. Since the set was very large, we used the  AFL function that minimizes inputs: \texttt{afl-cmin -i pngSuite/ -o seeds/ -- ./bin/png\_fuzz @@}. This command has reduced the number of images from $176$ to $80$.
For the dictionary, we used the one provided by AFL++.\endnote{AFL++ PNG dictionary: \texttt{AFLplusplus/dictionaries/png.dict}}
Next, in the Figure \ref{fig:status_fork}, one can see the AFL++ status screen, we see that the dictionary was used $612\;\mathrm{k}$ times and led the discovery of $9$ new paths when adding elements. However, it did not find any new paths by inserting words. AFL++ also did not find any new words to add to the dictionary. With the dictionary, we find $15$ new paths per million trial. In comparisons, with ‘avoc’, it found $1472$ path with $13.8\mathrm{M}$ trials. This gives $107$ discoveries per million trials. No splice have been used by the fuzzer.\\\\
The dictionary entries are ‘tokens’ or words specific to the domain of the library we are fuzzing (For example for libpng, one can think of the magic word/signature that must be present at the beginning of all PNGs (\texttt{89 50 4E 47 0D 0A 1A 0A}, which spells PNG and adds a bunch of new lines in order to find converting bugs to prevent invalid image as soon as possible). One can also think of the chunks types (e.g. \texttt{IHDR}, \texttt{IDAT}, \texttt{PLTE}). This method increases the likelihood of triggering new behaviors. We save a lot of time compared to only bit inversions, because it would take a long time to produce “89504E470D0A1A0A” in a completely random manner.